\chap The Back-Propagation Model

This short chapter aims to describe, in a step-by-step way,
the mathematical steps behind an implementation of the
back-propagation model.
The maths in the first sections is based upon Rumelhart et al., [20].
My interpretation of the maths, the proposed method of implementation,
appears in the last section of this chapter.
As in the last chapter, a biological analogy is drawn and processing
elements are referred to as \neus.

\sec The Maths Behind the Model

Based upon figure~7, which shows a feedforward \nn, the following 
mathematics holds:
\figure {7} {5} {A feedforward \nn}

Consider two neurons,~$i$ and~$j$, with \neu~$i$ being a lower layer
(that is closer to the input) than \neu~$j$.
Let the inputs to a \neu\ $j$ be $x_j$, and the outputs of a \neu~$i$ 
be $y_i$.
Let the weight of the connection between $i$ and $j$ be $w_{ji}$.
Then the total input $x_j$ into a \neu\ $j$ is a linear function of the
outputs $y_i$ of the units connected to $j$ and of the weights $w_{ji}$:
$$ x_j = \sum_i y_i w_{ji}.\eqno(1) $$
\noindent
Extra input or ``bias'' can be added to a \neu, say~$j$, equivalent to a
threshold of the opposite sign of its weight, assuming the input is~1.
The bias is treated just like any other $y$, ie $y_{\hbox{bias}} = 1,$
but $w_{j,\hbox{bias}}$ gives the negative threshold for $j$.

Each \neu\ $j$ has a non-integral output $y_j,$ which is a non-linear
function of its input $x_j:$
$$ y_j = {1 \over 1 + \hbox{e}^{-x_j}}.\eqno(2) $$
\noindent
The input/output function doesn't need to be the sigmoid response, any
non-linear function with a bounded derivative will do.

Let $c$ be an index of cases, $j$ be an index of output \neus, 
$y$ be the actual state of
an output \neu, and let $d$ be the desired state (or expected answer).
The total error, $E,$ is defined as
$$
E = {1 \over 2} \sum_c \sum_j (y_{j,c} - d_{j,c})^2. \eqno(3)
$$

The whole aim of this technique is to minimize the error, ie.\ make the
expected and actual outputs more similar.
To minimize $E$ by gradient descent, we need the partial derivative of
$E$ with respect to each weight in the network.
So from (3),
$$
{\partial E \over \partial y_j} = y_j - d_j. \eqno(4)
$$
And
$$
{\partial E \over \partial x_j} = {\partial E \over \partial y_j}
	{\partial y_j \over \partial x_j}. \eqno(5)
$$
From (2),
$$
{\partial y_j \over \partial x_j} = y_j(1-y_j).\eqno(6)
$$
Hence
$$
{\partial E \over \partial x_j} = {\partial E \over \partial y_j}
	y_j(1-y_j).\eqno(7)
$$
So we can calculate how a change in input $x$ to an output \neu\ will
affect the error.
However, the input to an output layer \neu\ is a linear combination of
the outputs from the lower layers and their weights.
So we can compute what the effect on the error would be for changes in
lower states and weights:
$$
{\partial E \over \partial w_{ji}} = {\partial E \over \partial x_j}
	{\partial x_j \over \partial w_{ji}} = {\partial E \over
		\partial x_j} y_i.\eqno(8)
$$
Then
$$
{\partial E \over \partial x_j} {\partial x_j \over \partial y_i}
	= {\partial E \over \partial x_j} w_{ji},\eqno(9)
$$
and taking into account all connections from $i,$
$$
{\partial E \over \partial y_i} = \sum_j {\partial E \over \partial x_j}
w_{ji}.\eqno(10)
$$
The procedure outlined above can be repeated for every layer below the
output layer, computing
$ \partial E \over \partial w$ as we go.

This brings us to the philosophies that can be employed to modify the
weights.

\sec Modification of Weights

The simplest scheme of weight modification is to modify them as we go
for every input/output pair. (To change~$w$, modify by~$\Delta w =
{\partial E \over \partial w}.$)
This method doesn't require $\partial E \over \partial w$ to be stored
for each pass.

Another method, used by Rumelhart et al.\ [20] is to accumulate
$\partial E \over \partial w$ over all the input/output pairs before
changing the weights.
The simplest version of this method is to make
$$
\Delta w = - \varepsilon {\partial E \over \partial w}, \eqno(11)
$$
where~$\varepsilon$ is a constant of proportionality and~$\partial E \over
\partial w$ has been accumulated over all cases.
An alternative version which apparently offers speed improvements is to
use a proportion of the previous~$\Delta w$,
$$
\Delta w(t) = -\varepsilon {\partial E \over \partial w(t)} + \alpha
			\Delta w(t-1), \eqno(12)
$$
\noindent
where~$t$ is a count of the number of times all input/output pairs have
been presented and~$\alpha$ is an exponential decay factor.
This is the method that I will use in the software model of the
back-propagation network.

\sec An Algorithm for The Back-Propagation Model

From all this mathematics, we can now explain the back-propagation model
algorithmically.
Two algorithms are presented.
I have found the first to cause networks to converge faster, but the
second may prove useful in some circumstances.
In the first algorithm~$\partial E \over \partial w$ is 
accumulated over all cases, and~$\Delta w(t)$ is calculated after
all input/output cases have been presented.
In the second algorithm,~$\partial E \over \partial w$ is
used to calculate~$\Delta w(t)$ after every input/output case.
\bigskip\noindent
The first algorithm is:
\smallskip
\item {1.} Set up random weights between \neus.
\item {2.} Set~$\sum {\partial E \over \partial w}$ to zero, and
set $E_{\rm TOTAL}$ to zero.
\item {3.} Input data sample and compare output with expected answer.
For all output \neus, if output differs from the expected value by
more than~$0.2$, increment $E_{\rm TOTAL}$.
\item {4.} If dealing with the weights for the top layer,
calculate~$\partial E \over \partial y$ from~$(4)$. Otherwise,
calculate~$\partial E \over \partial y$ from~$(10)$.
\item {5.} Calculate~$\partial E \over \partial x$ from~$(7)$.
\item {6.} Calculate~$\partial E \over \partial w$ from~$(8)$ and
add to~$\sum {\partial E \over \partial w}$.
\item {7.} If there are more input/output data pairs, go to~$3.$
again.
\item {8.} Calculate~$\Delta w(t)$ from~$(12)$, and save~$\Delta
w(t)$ as~$\Delta w(t-1)$ for next pass.
\item {9.} Apply~$\Delta w(t)$ to change weights,~$w$.
\item {10.} If~$E_{\rm TOTAL} = 0$, we are finished and the network
has ``learnt''.
Otherwise, go back to~$2.$ and repeat the procedure for the full set
of data pairs.
\bigskip\noindent
The second (alternative) algorithm is:
\smallskip
\item {1.} Set up random weights between \neus.
\item {2.} Set~$E_{\rm TOTAL}$ to zero.
\item {3.} Input data sample and compare output with expected answer.
For all output \neus, if output differs from the expected value by
more than~$0.2$, increment~$E_{\rm TOTAL}$.
\item {4.} If dealing with the weights for the top layer,
calculate~$\partial E \over \partial y$ from~$(4)$. Otherwise,
calculate~$\partial E \over \partial y$ from~$(10)$.
\item {5.} Calculate~$\partial E \over \partial x$ from~$(7)$.
\item {6.} Calculate~$\partial E \over \partial w$ from~$(8)$.
\item {7.} Calculate~$\Delta w(t)$ from~$(12)$, and save~$\Delta
w(t)$ as~$\Delta w(t-1)$ for next pass.
\item {8.} Apply~$\Delta w(t)$ to change weights,~$w$.
\item {9.} If there are more input/output data pairs, go to~$3.$
again.
\item {10.} If~$E_{\rm TOTAL} = 0$, we are finished and the network
has ``learnt''.
Otherwise, go back to~$2.$ and repeat the procedure for the full set
of data pairs.
\smallskip
The next part, Software, deals with the implemented software for the
back-propagation algorithm, where both the algorithms above have
been implemented.
